% part: intuitionistic-logic
% chapter: introduction
% section: bhk-interpretation

\documentclass[../../../include/open-logic-section]{subfiles}

\begin{document}

\olfileid[zh]{int}{int}{bhk}

\olsection{\zhFirst{Brouwer--Heyting--Kolmogorov}{布劳威尔–海廷–柯尔莫哥洛夫}解释}

\begin{editorial}
  使用 BHK 解释证明直觉主义命题的有效性令人困惑；必须把它们解释得更好。
\end{editorial}

对于直觉主义联结词，存在一种非形式的构造性解释，通常称为\zhFirst{Brouwer--Heyting--Kolmogorov}{布劳威尔–海廷–柯尔莫哥洛夫}解释，即 BHK 解释。它运用「构造」这一概念，你可以把它想成是一种构造性证明。（在 BHK 解释中我们不使用「证明」一词，以免与形式!!{derivation}系统中!!a{derivation}的概念相混淆。）基于这一直观概念，BHK 解释说明了直觉主义联结词的含义。

\begin{enumerate}
\item 我们假定，我们知道一个原子公式的构造由什么构成。
\item $!A_1 \land !A_2$ 的一个构造是一个二元组 $\tuple{M_1, M_2}$，其中 $M_1$ 是~$!A_1$ 的一个构造，$M_2$ 是~$!A_2$ 的一个构造。
\item $!A_1 \lor !A_2$ 的一个构造是一个二元组 $\tuple{s, M}$，其中 $s$ 为~$1$ 且 $M$ 是~$!A_1$ 的一个构造，或 $s$ 为~$2$ 且 $M$ 是~$!A_2$ 的一个构造。
\item $!A \lif !B$ 的一个构造是把~$!A$ 的一个构造变换成~$!B$ 的一个构造的函数。
\item 对于~$\lfalse$（荒谬）不存在构造。
\item $\lnot !A$ 定义为 $!A \lif \lfalse$ 的同义词。即，$\lnot !A$ 的一个构造是把~$!A$ 的一个构造变换成~$\lfalse$ 的一个构造的函数。
\end{enumerate}

\begin{ex}
以 $\lnot \lfalse$ 为例。它的一个构造是一个函数，它把 $\lfalse$ 的任意一个构造作为输入，输出的则是~$\lfalse$ 的一个构造。显然，恒等函数~$\Id{}$ 就是这样的构造：给定~$\lfalse$ 的一个构造~$M$，$\Id{}(M) = M$ 便得到一个~$\lfalse$ 的构造。
\end{ex}

一般说来，$\lnot !A$ 意为「$!A$ 的一个构造是不可能的」。

\begin{ex}
我们来证明对任意命题~$!A$ 有 $!A \lif \lnot \lnot !A$，即 $!A \lif ((!A \lif \lfalse) \lif \lfalse)$。该构造应是一个函数~$f$，它对 $!A$ 的一个构造~$M$ 输入，返回 $(!A \lif \lfalse) \lif \lfalse$ 的一个构造~$f(M)$。~$f$ 如下构造 $(!A \lif \lfalse) \lif \lfalse$ 的构造：我们必须定义一个函数 $g$，它把 $!A \lif \lfalse$ 的一个构造~$h$ 作为输入，输出~$\lfalse$ 的一个构造。我们可以如下定义 $g$：把输入~$h$ 作用于 $!A$ 的构造~$M$（即我们先前收到的那一个）。由于 $h$ 的输出~$h(M)$ 是 $\lfalse$ 的一个构造，因此若 $M$ 是~$!A$ 的一个构造，$f(M)(h) = h(M)$ 就是~$\lfalse$ 的一个构造。
\end{ex}

\begin{ex}
让我们为 $\lnot(!A \land \lnot !A)$，即 $(!A \land (!A \lif \lfalse)) \lif \lfalse$ 给出一个构造。这是一个函数~$f$，它把 $!A \land (!A \lif \lfalse)$ 的一个构造~$M$ 作为输入，产生~$\lfalse$ 的一个构造。一个合取式 $!B_1 \land !B_2$ 的构造是一个二元组 $\tuple{N_1, N_2}$，其中 $N_1$ 是~$!B_1$ 的一个构造，$N_2$ 是~$!B_2$ 的一个构造。我们可定义函数 $p_1$ 和 $p_2$，它们分别从 $!B_1 \land !B_2$ 的一个构造中恢复出 $!B_1$ 和 $!B_2$ 的构造：
\begin{align*}
  p_1(\tuple{N_1, N_2}) &= N_1\\
  p_2(\tuple{N_1, N_2}) &= N_2
\end{align*}
$f$ 做的事是这样的：首先把 $p_1$ 作用于其输入~$M$。这产生~$!A$ 的一个构造。然后把 $p_2$ 作用于~$M$，产生~$!A \lif \lfalse$ 的一个构造。这一构造本身是一个函数~$p_2(M)$，它若以 $!A$ 的一个构造作为输入，则产生~$\lfalse$ 的一个构造。换言之，若我们把 $p_2(M)$ 作用于 $p_1(M)$，就得到~$\lfalse$ 的一个构造。因此，我们可以定义 $f(M) = p_2(M)(p_1(M))$。
\end{ex}

\begin{ex}
让我们给出 $((!A \land !B) \lif !C) \lif (!A \lif (!B \lif !C))$ 的构造，即一个函数 $f$，它把 $(!A \land !B) \lif !C$ 的一个构造~$g$ 变换成 $(!A \lif (!B \lif !C))$ 的一个构造。构造 $g$ 本身是一个函数（从 $!A \land !B$ 的构造到~$C$ 的构造）。而输出 $f(g)$ 是一个函数~$h_g$，它把 $!A$ 的构造映到（把~$!B$ 的构造映到~$!C$ 的构造的）函数。

好的，这令人困惑。我们必须构造某个函数 $h_g$，它将是对输入~$g$ 时 $f$ 的输出。~$h_g$ 的输入是~$!A$ 的一个构造~$M$。$h_g(M)$ 的输出应该是一个函数~$k_M$，它把~$!B$ 的构造~$N$ 映到~$!C$ 的构造。令 $k_{g,M}(N) = g(\tuple{M, N})$。回想 $\tuple{M, N}$ 是~$!A \land !B$ 的一个构造。因此 $k_{g,M}$ 是 $!B \lif !C$ 的一个构造：它把~$!B$ 的构造~$N$ 映到~$!C$ 的构造。现在令 $h_g(M) = k_{g,M}$。这是一个函数，它把~$!A$ 的构造~$M$ 映到 $!B \lif !C$ 的构造 $k_{g,M}$。现在令 $f(g) = h_g$。这是一个函数，它把 $(!A \land !B) \lif !C$ 的构造 $g$ 映到 $!A \lif (!B \lif !C)$ 的构造。呼！{}
\end{ex}

语句 $!A \lor \lnot !A$ 称为排中律。我们可以对某些特定的 $!A$ 证明它（例如，$\lfalse \lor \lnot \lfalse$），但不能一般地证明。这是因为直觉主义析取要求给出其中一个析取项的构造，而有些语句目前既不能被证明也不能被反驳（例如，哥德巴赫猜想）。不过，你也不能反驳排中律：即，$\lnot \lnot (!A \lor \lnot !A)$ 成立。

\begin{ex}
要证明 $\lnot \lnot (!A \lor \lnot !A)$，我们需要一个函数~$f$，它把 $\lnot(!A \lor \lnot !A)$，即 $(!A \lor (!A \lif \lfalse)) \lif \lfalse$ 的一个构造，变换成~$\lfalse$ 的一个构造。换言之，我们需要一个函数 $f$，使得若 $g$ 是 $\lnot(!A \lor \lnot !A)$ 的一个构造，则 $f(g)$ 是~$\lfalse$ 的一个构造。

设 $g$ 是 $\lnot(!A \lor \lnot !A)$ 的一个构造，即一个把 $!A \lor \lnot !A$ 的一个构造变换成~$\lfalse$ 的一个构造的函数。$!A \lor \lnot !A$ 的一个构造是一个二元组 $\tuple{s, M}$，其中或者 $s = 1$ 且 $M$ 是 $!A$ 的一个构造，或者 $s = 2$ 且 $M$ 是 $\lnot !A$ 的一个构造。令 $h_1$ 为把 $!A$ 的一个构造~$M_1$ 映到 $!A \lor \lnot !A$ 的一个构造的函数：它把 $M_1$ 映到 $\tuple{1, M_1}$。令 $h_2$ 为把 $\lnot !A$ 的一个构造~$M_2$ 映到 $!A \lor \lnot !A$ 的一个构造的函数：它把 $M_2$ 映到 $\tuple{2, M_2}$。

令 $k$ 为 $\comp{h_1}{g}$：这是一个函数，它若给定~$!A$ 的一个构造，则返回~$\lfalse$ 的一个构造，即它是~$!A \lif \lfalse$ 或 $\lnot !A$ 的一个构造。现在令 $l$ 为 $\comp{h_2}{g}$。这是一个函数，它给定 $\lnot !A$ 的一个构造，则提供~$\lfalse$ 的一个构造。由于 $k$ 是 $\lnot !A$ 的一个构造，$l(k)$ 是~$\lfalse$ 的一个构造。

合起来，我们所做的是描述如何把 $\lnot(!A \lor \lnot !A)$ 的一个构造~$g$ 变换成~$\lfalse$ 的一个构造，即把 $\lnot(!A \lor \lnot !A)$ 的一个构造~$g$ 映到~$\lfalse$ 的构造 $l(k)$ 的函数 $f$，就是 $\lnot\lnot(!A \lor \lnot !A)$ 的一个构造。
\end{ex}

如你所见，用 BHK 解释来证明!!{formula}的直觉主义有效性很快就变得繁琐且令人困惑。幸运的是，对于直觉主义逻辑存在更好的!!{derivation}系统，也有更精确的语义解释。
\end{document}
